% ============================================================
% 分片：s01-intro  对应原文 PDF 第 3--13 页（印刷页 3--13）
% 负责章节：§1 Introduction + 1.1--1.9（§1 标题起）
% 译者：Lead   审校状态：初稿（样板分片）
% ============================================================

% ==== 原文 PDF 第 3 页（印刷页 3）====

\section{引言}

\begin{keytakeaway}
\textbf{核心物理图像：}几何 Langlands 对应可解释为两个对偶约化群 $G$ 与 ${}^{L}G$ 的 Hitchin 纤维化之间的镜像对称；一般纤维上的 T-对偶给出 B-膜范畴与 A-膜范畴的等价。当 B-模型一侧的 Hitchin 纤维出现轨形奇点时，相应的局部系统具有有限自同构群，即经典意义下的内窥局部系统；其在 A-模型一侧的对偶是奇异 Hitchin 纤维上的分数 A-膜。
\textbf{数学工具：}Hitchin 纤维化与 Higgs 丛、SYZ 镜像对称、A/B-膜范畴、D-模与反常层、Hecke 本征层、经典 Langlands 对应与 L-包、Arthur--Selberg 迹公式及其稳定化、Ngô 关于基本引理的几何工作。
\textbf{主要结论：}本文提出几何内窥的三条预言：(1) 内窥局部系统所对应的 Hecke 本征层按自同构群 $\Gamma$ 的不可约表示分裂，重数为表示的维数；(2) 分裂出的层满足「分数 Hecke 性质」，并在函数域情形得到验证；(3) 在一般正则半单点处，$\Gamma$ 同构于 $\pi_0(\boldsymbol{P}_b)$ 的对偶群的子群。
\end{keytakeaway}

本文研究 Langlands 对偶与镜像对称研究中出现的一些自然问题。在数学方面，问题是描述自守表示理论中内窥（endoscopy）现象的一个几何类比；在物理方面，则是考察以光滑对偶环面为目标流形的超对称 sigma 模型的 T-对偶，当这些环面变为奇异时的极限。

这两类问题结果密切相关，多少有些出人意料。原因在于：按照 [KW, GW]，几何 Langlands 对应可以解释为两个对偶约化群 $G$ 与 ${}^{L}G$ 的 Hitchin 纤维化之间的镜像对称：
\[
\begin{array}{ccccc}
\mathcal{M}_H({}^{L}G) & & & & \mathcal{M}_H(G)\\[2pt]
& \searrow & & \swarrow & \\[2pt]
& & \boldsymbol{B} & &
\end{array}
\]
这里 $\mathcal{M}_H(G)$ 表示光滑 Riemann 面 $C$ 上 Higgs $G$-丛的模空间，$\boldsymbol{B}$ 是对应的两个对偶 Hitchin 纤维化的公共底 [Hi1, Hi2]。二者之间的镜像对称由逐纤维的 T-对偶实现，处于 Strominger--Yau--Zaslow 一般图景 [SYZ] 的框架之内。这一对偶也与对应于 $G$ 与 ${}^{L}G$ 的某些超对称四维规范理论的 S-对偶密切相关 [KW]。

% [Tutor Note: 图景点拨。原文这句话把三个层次压缩在一起：物理上，$N=4$ 超 Yang--Mills 沿 $C$ 紧化给出以 $\mathcal{M}_H(G)$ 为目标流形的二维 sigma 模型；数学上，$\mathcal{M}_H(G)$ 与 $\mathcal{M}_H({}^{L}G)$ 都是超凯勒流形，关于复结构 $I$ 全纯、关于 $J$（或 $K$）则纤维为 special Lagrangian 环面。SYZ 图景只保证一般纤维上的 T-对偶；纤维退化时如何延拓正是本文的问题。图中三角形只是说两个纤维化共享底 $\boldsymbol{B}$，并不是说镜像对称是逐点给出的。]

\subsection{奇异纤维的 T-对偶}

当 $G=SL_n$ 时，两个 Hitchin 纤维化的一般纤维是光滑对偶环面（可描述为谱曲线的广义 Prym 簇）；对这些光滑纤维，T-对偶已理解得相当清楚。特别地，它在 $\mathcal{M}_H({}^{L}G)$ 于 $b\in\boldsymbol{B}$ 处的纤维 ${}^{L}\boldsymbol{F}_b$ 上的 B-膜范畴与对偶纤维 $\boldsymbol{F}_b$ 上的 A-膜范畴（更确切地说，是它们的连通分支）之间建立等价。
% ==== 原文 PDF 第 4 页（印刷页 4）====

最简单的 B-膜是支撑在 ${}^{L}\boldsymbol{F}_b$ 各点上的点凝聚层。\footnote{这里 B-膜范畴应在复结构 $J$ 下考虑，在此复结构下 $\mathcal{M}_H({}^{L}G)$ 实现为平坦 ${}^{L}G$-丛的模空间。不过，点层在任何复结构下都是合法的 B-膜。}在上述范畴等价下，它们对应于 $\boldsymbol{F}_b$ 上秩一酉平坦丛给出的 A-膜。后者具有一个重要性质：它们是某些算子的本征膜；这些算子是四维规范理论的 't Hooft 线算子 [KW] 的二维影子，并与 $G$-丛上的 Hecke 对应密切相关。与此对偶的性质是：点 B-膜是所谓 Wilson 算子 [KW] 的本征膜；后者更容易证明。

一个自然的问题是：在 Hitchin 纤维化的奇异纤维处，T-对偶呈现什么形态？特别地，T-对偶把支撑在 $\mathcal{M}_H({}^{L}G)$ 的奇异点上的 B-膜映到哪里？本文考虑纤维奇点尽可能温和的情形，即轨形奇点。（例如，当 $G=SL_n$ 时，若谱曲线不可约且既约，则这是唯一的奇点类型。）结果这恰好就是 [N1, N2] 中所定义的「椭圆内窥」情形（见下文）。本文详细描述 $G=SL_2$ 时的情形，并说明如何把结果推广到其他群。

当 $G=SL_2$ 时，我们关心的 $\mathcal{M}_H({}^{L}G)$ 的奇异点，对应于曲线 $C$ 上约化到子群 $O_2\subset SO_3$ 的 ${}^{L}G=SO_3$ 局部系统（或 Higgs 丛）；如下文所述，这是椭圆内窥最简单的可能情形。这类局部系统的一般成员具有自同构群 $\mathbb{Z}_2$（即 $O_2$ 的中心），因此 $\mathcal{M}_H({}^{L}G)$ 中对应的点实际上是 $\mathbb{Z}_2$-轨形点。这意味着支撑在该点上的 B-膜范畴等价于 $\mathbb{Z}_2$ 的表示范畴 $\mathrm{Rep}(\mathbb{Z}_2)$，从而有两个不可约对象。因此我们预期对偶的 A-膜范畴也应有两个不可约对象。事实上我们证明：此时对偶 Hitchin 纤维有两个不可约分支，所求的 A-膜正是支撑在这两个分支上的所谓分数膜（fractional branes）。只有它们的和（或并）才是 't Hooft--Hecke 算子的本征膜（这反映了如下事实：B-模型中 Wilson 算子的唯一本征膜对应于 $\mathbb{Z}_2$ 的正则表示，即其两个不可约表示的直和）。然而我们证明，两个分数 A-膜各自都满足标准 Hecke 性质的一种自然修正，即「分数 Hecke 性质」；对其他群它有直接的推广，并且具有独立的意义。

% [Tutor Note: 易错预警。这里出现本文的核心机制，务必分清两个范畴的方向：B-模型一侧的 $\mathbb{Z}_2$-轨形点，其 B-膜范畴是 $\mathrm{Rep}(\mathbb{Z}_2)$（两个不可约对象，分别对应两个分数膜）；A-模型一侧是 $\mathcal{M}_H(G)$ 的奇异 Hitchin 纤维，其对偶有两个不可约分支（在 $SL_2$、一个分歧点的情形下正是本文 §3 显式算出的两个分支）。「只有直和是本征膜」不是技术细节：Wilson 算子的本征膜对应正则表示，而正则表示分解为两个一维不可约表示，分数膜分别承载其一。作者此处把「同调镜像对称的范畴等价」与「物理膜的谱」混同叙述，阅读时须以 §3 的显式计算为准。]

\subsection{A-膜与 D-模}

在几何 Langlands 对应的通常表述中（例如见 [F1] 第 6 节），对应于 $C$ 上 ${}^{L}G$-局部系统的对象是所谓的 Hecke 本征层。它们是 $\Bun_G$（$C$ 上 $G$-丛的模叠）上满足 Hecke 性质的 D-模（或反常层）。然而其结构以复杂著称，难以显式分析。相反，在 [KW, GW] 发展的新形式体系中，Hecke 本征层被 't Hooft 算子的本征膜（即本征 A-膜）所取代。这些 A-膜更容易「实验观测」，也更容易显式分析。我们的想法是借助这一新语言来洞察几何 Langlands 对应的结构——特别是与内窥有关的部分（更雄心地说，是一般的函子性原理）。

从 $\mathcal{M}_H(G)$ 上的 A-膜到 $\Bun_G$ 上的 D-模的过渡已在 [KW] 中解释（另见 [NZ, Nad] 中可能的替代进路；也见 [ADKMV, DHSV]，那里从另一个角度把 D-模引入物理）。虽然这一过渡尚未在数学上完全严格化，但它已足以描述与本征膜相联系的 Hecke 本征层的重要特征，例如其可约性、$\Bun_G$ 上 Hecke 本征层由局部系统表示的开子集、这些局部系统的秩，乃至其单值。于是我们关于 A-膜的结论对 Hecke 本征层有直接推论。特别地，若一个本征膜 $\mathcal{A}$ 分解为两个不可约膜 $\mathcal{A}_1$ 与 $\mathcal{A}_2$，则我们预言相应的 Hecke 本征层 $\mathcal{F}$ 也分解为两个 D-模 $\mathcal{F}_1$ 与 $\mathcal{F}_2$ 的直和，分别对应于 $\mathcal{A}_1$ 与 $\mathcal{A}_2$。进一步，这两个 D-模应各自满足上面提到的分数 Hecke 本征层性质。

% [Tutor Note: 含糊点拨。原文说 A-膜到 D-模的过渡「尚未完全严格化」，被省略的是其物理来源：在 $\mathcal{M}_H(G)$ 上取典范余迷向膜 $\mathcal{A}_{cc}$，把 A-膜 $\mathcal{A}$ 映为 $\mathrm{Hom}(\mathcal{A}_{cc},\mathcal{A})$，沿 Hitchin 纤维化把开弦态空间层化到 $\Bun_G$，即得扭 D-模；$\mathcal{A}_{cc}$ 的自同态代数给出 $\Bun_G$ 上的整体微分算子。严格化需要半无限上同调与微局部层论，见 [NZ, Nad]。作者说「足以描述可约性、秩、单值」正是指这一层的函子性已足以支撑 §3--§6 的范畴层面论证。]

\subsection{从 \texorpdfstring{$\mathbb{C}$}{C} 上的曲线到 \texorpdfstring{$\mathbb{F}_q$}{Fq} 上的曲线}

上述一切的要旨是：通过分析支撑在奇异 Hitchin 纤维上的 A-膜范畴，我们得以洞察几何 Langlands 对应。随后我们再作一次信念的飞跃，假设我们在 $\mathbb{C}$ 上曲线情形观察到的 Hecke 本征层上的结构（例如分解为两个直和项）对有限域上的曲线也应成立。在后一情形，我们可通过取 Frobenius 在茎上的迹，把自守函数联系到分数 Hecke 本征层（这被称为 Grothendieck 的 faisceaux--fonctions 词典，见 §8.1）。于是关于分数 Hecke 本征层的预言转化为关于这些自守函数在经典 Hecke 算子作用下行为的具体预言。我们证明满足这些性质的函数确实存在，从而为我们的猜想提供了一致性检验。

因此，我们的出发点是两个对偶 Hitchin 纤维化上膜范畴之间的同调镜像对称。通过施加如下变换序列：
% ==== 原文 PDF 第 5 页（印刷页 5）====
\begin{equation}\tag{1.1}
\begin{array}{ccccccc}
\boxed{\text{A-branes}} & \xRightarrow{\ \text{over }\mathbb{C}} &
\boxed{\CD\text{-modules}} & \xRightarrow{\ \text{over }\mathbb{C}} &
\boxed{\begin{smallmatrix}\text{perverse}\\ \text{sheaves}\end{smallmatrix}} & \xRightarrow{\ \text{over }\mathbb{F}_q} &
\boxed{\begin{smallmatrix}\text{automorphic}\\ \text{functions}\end{smallmatrix}}
\end{array}
\end{equation}
我们就把在这一镜像对称研究中观察到的 A-膜结构与经典自守形式理论联系起来。从这个观点看，作为 't Hooft--Hecke 算子本征膜（通常意义下）的 A-膜，是概括不可约自守表示的那些 Hecke 本征函数的几何类比。

% [Tutor Note: 易错预警。链 (1.1) 中三个箭头的含义并不相同，不可当作逐项恒等：第一个箭头是 [KW] 的 A-膜/扭 D-模对应（部分仍是猜想）；第二个箭头是 Riemann--Hilbert 对应（$\mathbb{C}$ 上 D-模与反常层之间的范畴等价，依赖正则奇性条件）；第三个箭头才是 Grothendieck 的 faisceaux--fonctions 词典（$\mathbb{F}_q$ 上：取 Frobenius 在茎上的迹）。注意第二个箭头两侧基域相同（均为 $\mathbb{C}$），但第三个箭头换到了 $\mathbb{F}_q$，且只保留「函数」这一层信息，故链条末端丢失了范畴结构。作者用同一条链把几何内窥的预言搬到经典理论，风险也正在这里。]

由此引出一个诱人的问题：在内窥情形下，当 Hitchin 纤维变为奇异时，A-本征膜所分裂成的分数本征膜，其表示论类比是什么？R. Langlands 此前曾提出不可约自守表示与基本粒子之间的一个漂亮类比 [L3]。从这个观点看，分数 A-膜的存在暗示自守表示具有某种内部的、「类夸克」的结构；这一结构仍有待充分探索与理解。
% ==== 原文 PDF 第 6 页（印刷页 6）====

我们希望通过理解这些结构，获得关于内窥的几何意义的重要线索。

\subsection{经典内窥}

内窥是经典 Langlands 对应中最引人入胜的现象之一。为说明其动机，我们回顾（详见例如 [F1] 第 2.4 节）群 $GL_n$ 的 Langlands 对应。设 $C$ 是有限域 $\mathbb{F}_q$ 上的光滑射影曲线，$F$ 是 $C$ 上的有理函数域。Weil 群 $W_F$ 是 $F$ 的（可分）闭包 $\bar F$ 的自同构群 $\mathrm{Gal}(\bar F/F)$ 的稠密子群。Langlands 对应在 $W_F$ 的 $n$ 维（$\ell$-进）表示与 $GL_n(\mathbb{A}_F)$ 的不可约自守表示之间建立双射，其中 $\mathbb{A}_F$ 是 $F$ 的 adele 环 [Dr1, Dr2, Laf]。\footnote{这是函数域的 Langlands 对应。数域情形有类似但更复杂的版本，其中 $F$ 换成有理数域 $\mathbb{Q}$ 或其有限扩张；详见例如 [F1] 第一部分。}若把 $GL_n$ 换成更一般的约化群 $G$，则在一级近似下，我们应预期 $G(\mathbb{A})$ 的不可约自守表示与从 $W_F$ 到 $G$ 的 Langlands 对偶群 ${}^{L}G$ 的（$\ell$-进）同态 $\sigma$ 双射。然而事实上，一般而言每个 $\sigma$ 对应的不是一个、而是若干个（可能无穷多个）$G(\mathbb{A}_F)$ 的不可约自守表示。这些表示的等价类之集称为对应于 $\sigma$ 的 L-包；这一名称来自 Labesse--Langlands [LL] 的工作，他们发现了这一现象（对群 $G=SL_2$）。

L-包的结构在如下情形最为有趣：同态 $W_F\to{}^{L}G=SO_3$ 的像包含在子群 $O_2\subset PGL_2$ 中，但不包含在其连通分支 $SO_2$ 中；这意味着其自同构群不连通（一般为 $\mathbb{Z}_2$）。于是 $SL_2(\mathbb{A}_F)$ 的自守表示理论部分地由群 $O_2$ 支配。这个群（或者按传统术语，其对偶为 $O_2$ 的 $SL_2(F)$ 的子群 $H$）是内窥群的一个例子，而这一关系就是被称为「内窥」的神秘现象的一个例子。它由 Langlands 等人在试图以与迹公式几何侧出现的轨道积分结构相容的方式组织自守表示时发现。下面简要解释这一点，细节请读者参阅 [L2, Ko, Art2]。迹公式，或者更确切地说，它的「正则椭圆部分」（本文仅讨论这一部分），具有如下一般形式：

% [Tutor Note: 跳步点拨。原文「一级近似下应为双射，但实际上每个 $\sigma$ 对应若干表示」正是经典内窥的入口：从 $W_F$ 到 ${}^{L}G$ 的同态只固定了 L-包的整体参数，包内各成员由内窥群（即 ${}^{L}G$ 中半单元素的中心化子）区分。对 $SL_2$ 的情形，$L$-包由 Labesse--Langlands 首次显式算出；其重数由 Kottwitz 猜想的重数公式给出（见 (1.4) 后的讨论）。]

% ==== 原文 PDF 第 7 页（印刷页 7）====
\begin{equation}\tag{1.2}
\boxed{\begin{smallmatrix}\text{谱侧}\\ \text{(spectral side)}\end{smallmatrix}}
\;=\;
\boxed{\begin{smallmatrix}\text{几何侧}\\ \text{(geometric side)}\end{smallmatrix}}
\end{equation}
谱侧等于一个紧支撑检验函数 $f\in C_c^\infty(G(\mathbb{A}_F))$ 在 $G(\mathbb{A}_F)$ 的不可约缓增尖点自守表示上的迹之和（我们回忆，这些表示实现于商 $G(F)\backslash G(\mathbb{A}_F)$ 上的某个函数空间）：
\begin{equation}\tag{1.3}
\sum_{\sigma:\,W_F\to{}^{L}G}\ \sum_{\phi\in L_\sigma} m_\phi\,\operatorname{Tr}(f,\pi_\phi).
\end{equation}
这里求和遍历某一类同态 $\sigma:W_F\to{}^{L}G$，它们应当标记 L-包 $L_\sigma$，即不可约自守表示的等价类之集 $\{\pi_\phi\}_{\phi\in L_\sigma}$；$m_\phi$ 表示 $\pi_\phi$ 在自守函数空间中的重数。

几何侧是 $f$ 的轨道积分之和，即 $f$ 在 $G(F)$ 的元素于 $G(\mathbb{A}_F)$ 中的共轭类上的积分。

几何侧需要「稳定化」。这意味着把它改写为 $G(F)$ 的元素在 $G(\mathbb{A}_F)$ 中的稳定共轭类上的积分之和。\footnote{设 $k$ 为任意域。$G(k)$ 中两个元素称为稳定共轭的，如果它们在 $G(\bar k)$ 中共轭。由于 $\mathbb{A}_F$ 是 $F$ 的各完备化的限制乘积，我们同样得到 $\mathbb{A}_F$ 中稳定共轭的自然概念。}这在许多理由下都是必要的；理由之一是，不做这一步甚至无法期望比较不同群的迹公式的几何侧（例如见 [Art2] 第 27 节）。几何侧所得的表达式为 [L2]
\begin{equation}\tag{1.4}
\sum_{H}\iota(H,G)\,ST^{G\text{-}\mathrm{reg}}_{\mathrm{ell}}(f^H),
\end{equation}
其中求和遍历 $G$ 的椭圆内窥群 $H$（它们一般不是 $G$ 的子群），并包括 $H=G$ 本身；$ST^{G\text{-}\mathrm{reg}}_{\mathrm{ell}}(f^H)$ 表示群 $H$ 的稳定轨道积分之和，$\iota(H,G)$ 是某些数。椭圆内窥群粗略地定义为 ${}^{L}G$ 中半单元素的中心化子的对偶群，这些半单元素不包含在 ${}^{L}G$ 的任何真 Levi 子群中。

% [Tutor Note: 易错预警。(1.3) 与 (1.4) 是同一个迹公式的两侧，但不可逐项比较：谱侧按同态 $\sigma$（即 L-包）分组，几何侧按内窥群 $H$ 分组，两者之间由稳定化（transfer）联系。$\iota(H,G)$ 是 Langlands--Shelstad 转移因子；$ST$ 的下标 $\mathrm{ell}$ 指只取椭圆元素（正则椭圆部分），上标 $G\text{-}\mathrm{reg}$ 指 $G$-正则。注意 (1.4) 中的求和包含 $H=G$，而 §1.5 的几何内窥定义刻意排除 $G$（见原书 p.8 脚注 4），两者语境不同，不要混淆。]

公式 (1.4) 依赖于若干假设，其中最重要的是所谓的转移猜想。它断言存在一个映射 $f\mapsto f^H$，从 $G(\mathbb{A}_F)$ 上的函数到 $H(\mathbb{A}_F)$ 上的函数，使得粗略地说，$f^H$ 的稳定轨道积分等于 $f$ 的稳定轨道积分经某种扭转修正后的值（细节见 [L2, Art2, Dat] 及其参考文献）。

转移猜想的一个特殊情形（当 $f$ 是 $G$ 的极大紧子群的特征函数时；此时要求 $f^H$ 也是同类函数）就是所谓的基本引理。基本引理（在函数域情形）最近由 B.-C. Ngô [N2] 证明（另见 [GKM, La5, LN, N1]）。更确切地说，Ngô 证明的是基本引理的 Lie 代数版本，而 Waldspurger 证明了它蕴含群版本以及一般的转移猜想 [Wa1]。此外，我们正在讨论的函数域情形的基本引理，在剩余域特征足够大时等价于数域情形的基本引理 [Wa2, CL]。
% ==== 原文 PDF 第 8 页（印刷页 8）====

由于迹公式的几何侧具有 (1.4) 的形式，自然预期谱侧也有类似形式，即可写为以椭圆内窥群标记的项之和，且 (1.2) 两边的对应项相等。这类公式最先由 Labesse--Langlands 对 $G=SL_2$（及其内形式）建立。一般情形下的相应公式由 R. Kottwitz 猜想 [Ko]。该公式对自守表示理论有许多重要推论。首先，它给出 L-包中自守表示的重数 $m_\phi$ 的显式公式，其中 L-包对应于同态 $\sigma:W_F\to{}^{L}G$（出现在 (1.3) 中）。答案是：与那些满足 $\sigma$ 的像包含在 ${}^{L}H\subset{}^{L}G$ 中的椭圆内窥群 $H$ 相联系的项的线性组合。

这里或许有助于解释 $\sigma$ 与内窥群之间的关系。为简单起见，设 ${}^{L}G$ 是伴随型半单群（因而中心平凡），且 (1.3) 中出现的同态 $\sigma:W_F\to{}^{L}G$ 的像具有有限中心化子 $S_\sigma$。则与 $\sigma$ 相联系的各内窥群 $H$ 的 Langlands 对偶 ${}^{L}H$ 恰好是 ${}^{L}G$ 中 $S_\sigma$ 的非平凡元素的中心化子（因此 $\sigma$ 的像自动包含在 ${}^{L}H$ 中）。例如，若 ${}^{L}G=SO_3$，则这样得到的唯一子群是 $O_2\subset SO_3$。\footnote{注意，与 [N1, N2] 相同而与标准约定相反，我们在此不把 $G$ 本身视为椭圆内窥群。}

转移猜想与稳定化迹公式的第二个、或许更重要的推论，是它给出了 Langlands 函子性原理 [L1] 对同态 ${}^{L}H\to{}^{L}G$ 的具体实现，其中 $H$ 取遍椭圆内窥群。也就是说，我们得到从 $H(\mathbb{A}_F)$ 的自守表示的 L-包到 $G(\mathbb{A}_F)$ 的自守表示的 L-包的自然映射（详见例如 [Art1] 与 [Art2] 第 26 节）。

总之，经典内窥在 $G(\mathbb{A}_F)$ 的自守表示与其内窥群 $H(\mathbb{A}_F)$ 的自守表示之间建立了一种难以捉摸的联系；通过迹公式，这一联系与转移猜想所提供的两群轨道积分之间的关系相匹配。

\subsection{几何内窥}

过去二十年中，把经典 Langlands 对应翻译成几何语言取得了重大进展。正在形成的几何 Langlands 对应的优点是，它不仅对有限域上定义的曲线有意义，而且对复数域上的曲线（即 Riemann 面）也有意义。在这一版本中，我们有机会运用复代数几何的丰富资源，从而推进对一般 Langlands 对偶图景的理解。

经典 Langlands 对应的许多概念已在几何版本中找到对应物。一个尚未在几何上被理解的重要现象就是内窥。本文的目标是在发展几何内窥方面迈出第一步；所谓几何内窥，我们指的是揭示内窥群在几何 Langlands 对应中所扮演的特殊角色。由于几何情形中没有迹公式的明显类比，我们并不试图模仿导致经典内窥的迹公式稳定化。相反，我们希望描述在椭圆内窥型 ${}^{L}G$-局部系统的几何 Langlands 对应中出现的新结构。当 ${}^{L}G$ 半单时，这意味着（与经典情形类比）该局部系统的自同构群是严格大于 ${}^{L}G$ 的中心的有限群。例如，当 ${}^{L}G=SO_3$ 时，这些是像包含在 $O_2\subset SO_3$ 中但不包含在 $O_2$ 的极大环面中的局部系统。此时自同构群等于 $O_2$ 的中心 $\mathbb{Z}_2$，除非像包含在子群 $\mathbb{Z}_2\times\mathbb{Z}_2$ 中（本文基本忽略这一情形）。
% ==== 原文 PDF 第 9 页（印刷页 9）====

我们的方法是利用与 $G$ 及 ${}^{L}G$ 相联系的 Hitchin 纤维化的镜像对称，探索对应于内窥 ${}^{L}G$-局部系统的 A-膜结构；这些局部系统在 B-模型一侧被视为 $\mathcal{M}_H({}^{L}G)$ 的轨形点，在 A-模型一侧则实现于 $\mathcal{M}_H(G)$ 中相应的奇异 Hitchin 纤维上。这一方法的第一个优点是，内窥群（在经典理论中相当神秘，它们是在对群 $G$ 的迹公式作稳定化的过程中出现的）在此是显然可见的：它们自然地出现在镜像对称的 B-模型一侧。事实上，对任何子群 ${}^{L}H\subset{}^{L}G$，都有自然嵌入 $\mathcal{M}_H({}^{L}H)\hookrightarrow\mathcal{M}_H({}^{L}G)$；而对应于内窥局部系统 $E$ 的 $\mathcal{M}_H({}^{L}G)$ 中的点，属于所有与它相联系的内窥群 $H$（即那些 ${}^{L}H$ 包含 $E$ 的像的 $H$）的 $\mathcal{M}_H({}^{L}H)$ 的像。第二个优点（上面已提到）是，至少在一股正则半单情形，相应的 A-膜具有简单透明的结构（与 Hecke 本征层相反），这大大简化了我们的分析。

% [Tutor Note: 背景补充。这里「内窥群在 B-模型一侧显然可见」是本文的核心观察：经典内窥的稳定化需要人为选取转移因子 $f\mapsto f^H$，而在几何图景中，$H$ 直接表现为 $\mathcal{M}_H({}^{L}G)$ 的一个子流形 $\mathcal{M}_H({}^{L}H)$——即那些自同构群包含 $H$ 的局部系统的轨道。换言之，内窥不是附加的簿记，而是 Hitchin 模空间分层（按自同构群分层）的几何。]

随后我们把在 A-膜范畴（A-模型一侧）中观察到的结构，用 $\Bun_G$ 上的 D-模（对 $\mathbb{C}$ 上的曲线）或反常层（对 $\mathbb{C}$ 上或 $\mathbb{F}_q$ 上的曲线）来解释；它们是几何 Langlands 对应中更标准的对象。\footnote{或者，也可以从 A. Beilinson 与 V. Drinfeld 所建议的 Fourier--Mukai 变换的非阿贝尔版本的角度来看待它 [La4, Ro]。}最后，对于 $\mathbb{F}_q$ 上的曲线，我们研究联系到这些反常层的自守函数（见示意图 (1.1)）。由此得到一系列具体预言：

\begin{itemize}[leftmargin=2em,itemsep=2pt,topsep=2pt]
\item \textbf{第一个预言：}对应于自同构群为有限群 $\Gamma$ 的椭圆内窥 ${}^{L}G$-局部系统 $E$ 的 Hecke 本征层，分裂为以 $\Gamma$ 的不可约表示 $R$ 标记的不可约层 $\mathcal{F}_R$ 的直和，且 $\mathcal{F}_R$ 的重数等于 $\dim R$。
\item \textbf{第二个预言：}层 $\mathcal{F}_R$ 满足第 6 节所述的分数 Hecke 性质。这一预想在 $\mathbb{F}_q$ 上的曲线情形得到证实：我们验证了：赋值给这些层、且满足分数 Hecke 性质之函数论类比的那些函数确实存在。此外，我们把它们表示为通常 Hecke 本征函数经由某种关于 $\Gamma$ 的 Fourier 变换得到的线性组合。\footnote{这多少让人想起 G. Lusztig 在特征层理论中观察到的 Fourier 变换 [Lu]。}于是，在内窥情形，赋予不可约反常层的函数并非 Hecke 本征函数，而是其线性组合。
% ==== 原文 PDF 第 10 页（印刷页 10）====
\item \textbf{第三个预言：}若 $E$ 在 Hitchin 底 $\boldsymbol{B}$ 中的像 $b$ 是一般正则半单点，则群 $\Gamma$ 是阿贝尔的，且同构于 $\pi_0(\boldsymbol{P}_b)$ 的对偶群的子群，其中 $\boldsymbol{P}_b$ 是 $b$ 的广义 Prym 簇。此外，对偶 Hitchin 纤维 ${}^{L}\boldsymbol{F}_b$ 中存在一个局部系统，使 $\Gamma$ 同构于 $\pi_0(\boldsymbol{P}_b)$ 的对偶群（注意 $\pi_0(\boldsymbol{P}_b)$ 在 $\mathcal{M}_H(G)$ 中 $b$ 上的 Hitchin 纤维 $\boldsymbol{F}_b$ 的一个开稠子集上单可迁地作用，且该纤维在此情形是既约的 [N2]）。\footnote{在本文写作过程中，B.-C. Ngô 告知我们，他也已知道这一陈述。}
\end{itemize}

我们希望对这些现象的正确理解能导致对内窥以及基本引理等相关课题的更好理解。

\subsection{与 B.-C. Ngô 工作的联系}

我们的分析与经典内窥之间的联系来自如下事实：我们所用的几何与 Ngô Bảo Châu（吴宝珠）近期工作 [N1, N2] 中所利用的几何相似（另见优秀综述 [Dat]）。Ngô 发现了迹公式几何侧出现的轨道积分与模空间 $\mathcal{M}_H(G)$ 中 Hitchin 纤维的上同调之间引人注目的联系（更确切地说，是在 $\mathcal{M}_H(G)$ 的推广版本中，这些版本参数化极点除子 $D$ 足够大的亚纯 Higgs 场）。他借此在函数域情形、Lie 代数设定下证明了基本引理（对酉群，他此前与 G. Laumon 合作完成 [LN]）。他的做法是把轨道积分解释为有限域上曲线定义的 Higgs 丛模叠 $\mathcal{M}_H(G)$ 中 Hitchin 纤维 $\boldsymbol{F}_b$ 的点数；这些点数又解释为 Frobenius 作用在 $\boldsymbol{F}_b$ 的（$\ell$-进）上同调上的迹。

Ngô 构造中的关键一步，是按照有限阿贝尔群 $\pi_0(\boldsymbol{P}_b)$ 的作用分解该上同调，其中 $\boldsymbol{P}_b$ 是与椭圆点 $b\in\boldsymbol{B}$ 相联系的广义 Prym 簇。他把分解中 $\kappa$-同型部分（$\kappa$ 是 $\pi_0(\boldsymbol{P}_b)$ 的特征标）与内窥群 $H_\kappa$ 的相应 Hitchin 纤维上同调分解的 $\kappa$-部分等同起来。取 Frobenius 在这些子空间上的迹，就得到基本引理。（这里应提到更早的工作 [GKM, La5]，其中已给出基本引理密切相关的几何解释。）

% [Tutor Note: 跳步点拨。这里省略的是 Ngô 基本引理证明的完整逻辑链：轨道积分 = 有限域上 Hitchin 纤维的点数（Borisov--Libgober 型计数）$\to$ 由 Lefschetz 迹公式 = Frobenius 在上同调上的迹 $\to$ 按 $\pi_0(\boldsymbol{P}_b)$ 的 $\kappa$-同调分解（Ngô 的支撑定理/几何稳定化）$\to$ $\kappa$-部分与内窥群 $H_\kappa$ 的 Hitchin 纤维上同调的同构（这是 Ngô 定理的核心）$\to$ 逐项比较得基本引理。本文只需要这一几何事实的「形状」，不需要其证明细节。]

因此，Ngô 所用的几何看来与我们所使用的几何非常接近。事实上，我们考虑 $\mathcal{M}_H(G)$ 上支撑在 Hitchin 纤维 $\boldsymbol{F}_b$ 上的分数 A-膜，它们本质上以 $\pi_0(\boldsymbol{P}_b)$ 标记；而 Ngô 考虑类似定义的 Hitchin 纤维的上同调及其在 $\pi_0(\boldsymbol{P}_b)$ 下的分解。然而存在重要差别。首先，我们在 $\mathbb{C}$ 上工作，而 Ngô 在 $\mathbb{F}_q$ 上工作；在后一设定中，$\mathcal{M}_H(G)$ 与 $\mathcal{M}_H({}^{L}G)$ 之间的同调镜像对称没有明显的类比，而这一镜像对称对我们的方法至关重要。（事实上，对偶群 ${}^{L}G$ 与对偶 Hitchin 模空间 $\mathcal{M}_H({}^{L}G)$ 在 Ngô 的工作中都不起作用。）其次，Ngô 使用以曲线 $C$ 上有效除子 $D$ 标记的广义 Higgs 丛模空间；我们的模空间 $\mathcal{M}_H(G)$ 对应于 $C$ 上典范线丛的一个除子。\footnote{严格说来，这一情形超出 Ngô 工作的范围，因为他要求 $\deg(D)>2g-2$。但据他向我们解释，当 $D$ 取典范线丛的除子时，他的大部分结果仍然成立。}第三，也最重要，我们赋予奇异 Hitchin 纤维 $\boldsymbol{F}_a$ 各连通分支的对象——A-膜——是自守性质的对象；我们希望把它们与 Hecke 本征层、并最终与经典理论中的自守函数联系起来。因此这些对象应位于迹公式 (1.2) 的谱侧。另一方面，Ngô 把 Hitchin 纤维的点数（及其上同调）与迹公式几何侧（更确切地说其 Lie 代数版本）出现的轨道积分联系起来。先验地，这与自守表示（或自守函数）毫无关系！轨道积分与自守表示之间的联系由迹公式提供，但方式是相当间接的组合方式；只有当我们在几何侧对所有轨道积分求和、在谱侧对所有表示求和时，它才出现。
% ==== 原文 PDF 第 11 页（印刷页 11）====

这就提出如下问题：$\mathbb{F}_q$ 上 $\mathcal{M}_H(G)$ 中的单个 Hitchin 纤维与单个自守表示之间是否可能有直接联系？毕竟，我们有从 $\mathcal{M}_H(G)$ 到 $\Bun_G$ 的自然遗忘映射，而未分歧自守函数正生活在 $\Bun_G$ 上。上面讨论的从 A-膜到 Hecke 本征层的过渡，在经典理论中是否可能有类比，即从轨道积分到 Hecke 本征函数的过渡？

无论如何，我们觉得值得注意的是：Ngô 用来理解迹公式几何侧的同一个 Hitchin 纤维化几何，也被我们用于通过镜像对称研究几何内窥，因而（经由对应 (1.1)）似乎与自守表示相关。这一联系可能意义重大，因为它可能为迹公式以及一般的自守表示理论提供新的启示。\footnote{这似乎与 R. Langlands 在他最近的邵逸夫奖演讲 [L4] 中所表达的观点相呼应。}

\subsection{量子场论}

最后，我们希望把计算与量子场论中的各种问题联系起来。回忆 Strominger--Yau--Zaslow 图景 [SYZ] 把两个流形 $X$ 与 $Y$ 的同调镜像对称与 $X$、$Y$ 中对偶 special Lagrangian 纤维化的 T-对偶联系起来。这对一般纤维（它们是光滑的）尤其有效。这些纤维的 T-对偶可视为镜像对称的一种阿贝尔版本（与 Fourier--Mukai 变换密切相关）。因此，关于 $X$ 与 $Y$ 的镜像对称的重要「非阿贝尔」信息隐藏在奇异纤维的对偶之中。special Lagrangian 纤维化一般难以理解，但在 Hitchin 纤维化的情形，超凯勒结构导致极大的简化 [HT]，我们将利用它来分析某些奇异纤维的情形。特别地，我们证明在镜像对称下，$X$ 上的轨形奇点对应于 $Y$ 中可约的拉格朗日纤维；我们相信这是一个相当普遍的现象。从这个观点看，几何内窥表现为带轨形奇点的镜像对称的一个特殊情形。

% [Tutor Note: 易错预警。原文「$X$ 上的轨形奇点对应 $Y$ 中可约的拉格朗日纤维」是该文的一般性猜想，并非定理；能严格化的是 Hitchin 纤维化这个特例。要点在超凯勒旋转：$\mathcal{M}_H(G)$ 关于复结构 $I$ 是复辛流形（B-模型的靶），关于 $J$ 则其 Hitchin 纤维是 special Lagrangian 环面（A-模型的靶）；同一流形在不同复结构下分别是 $X$ 与 $Y$ 的角色。若忽略这一点，就会误以为「可约纤维」只是几何现象，而非镜像对偶的必然结果。]

几何内窥或许可以通过超对称畴壁获得自然的物理实现：畴壁一侧是规范群为 ${}^{L}H$ 的 $N=4$ 超对称 Yang--Mills 理论，另一侧是规范群为 ${}^{L}G$ 的同一理论（其中 $H$ 是 $G$ 的内窥群）。在对偶之下，该畴壁尤其应给出从 $\mathcal{M}_H(H)$ 上 A-膜范畴到 $\mathcal{M}_H(G)$ 上 A-膜范畴的函子。这是 Langlands 纲领中转移概念、即函子性原理的几何类比（见 §5.2.6）。

最后，正如我们在 §3.11 所解释的，本文研究的 Hitchin 纤维化也出现在 Seiberg--Witten 理论中。在该语境下，重要的四维流形不变量来自内窥点的贡献。同一 Hitchin 纤维化出现在两个不同问题中，其根源在于一个底层的六维量子场论，它可以以两种不同方式紧化到四维。
% ==== 原文 PDF 第 12 页（印刷页 12）====

\subsection{全文安排}

第 2 节概述对偶 Hitchin 纤维化的同调镜像对称与几何 Langlands 纲领之间的联系。第 3 节讨论我们的主要例子：椭圆曲线上带一个驯顺分歧点的 $SL_2$ 与 $SO_3$ Higgs 丛的模空间。我们详细描述这些模空间、其奇点以及相应的膜范畴。这个例子将作为本文所发展的普遍图景的原型。

第 4 节更详细地讨论从 A-膜到 D-模的过渡。虽然我们主要聚焦前一节发展的亏格一例子，但这一讨论的许多方面适用于更一般的设定。第 5 节把我们结果推广到更高亏格的曲线，并显式计算 Wilson 算子与 't Hooft/Hecke 算子对与几何内窥相关的电膜与磁膜的作用。第 6 节解释这些结果如何纳入一般的范畴形式体系；特别地，我们引入「分数 Hecke 本征层」的概念，并猜想第 5 节所发现的分数 A-膜对应的 D-模属于这一类。

接下来的任务是描述这些猜想在曲线 $C$ 定义于有限域时的类比，并把它们与自守函数理论联系起来。第 7 节回顾自守形式经典理论中内窥与 L-包的框架，特别聚焦于 Labesse 与 Langlands [LL] 最初分析的 $G=SL_2$ 非分歧情形。第 8 节讨论前面各节几何图景对经典理论的潜在推论，特别是计算与分数 Hecke 本征层相联系的自守函数。第 9 节把结果与猜想推广到其他群。最后，第 10 节讨论我们在 $SL_2$ 情形发现的两张分数本征膜不可区分这一微妙之处。我们把这一现象追溯到某个 $\mathbb{Z}_2$-gerbe，它看来是 Hitchin 纤维化镜像对称中一个微妙但重要的要素。

% ==== 原文 PDF 第 13 页（印刷页 13）====
在附录中，我们用 Whittaker 函数具体解释 $SL_2$ 的非分歧整体 L-包的结构。

\subsection{致谢}

E.F. 感谢 R. Langlands 向他传达关于内窥的洞见，也感谢 J. Arthur、D. Kazhdan、R. Kottwitz、B.-C. Ngô、M. Olsson、T. Pantev、B. Poonen 与 K. Ribet 对与本文相关的各种问题的有益回答，以及 D. Arinkin、A. Beilinson、R. Bezrukavnikov 与 V. Lafforgue 的启发性讨论。E.W. 同样感谢 T. Hausel、D. R. Morrison、T. Pantev、I. M. Singer 与 C. Vafa。我们也感谢 S. Gukov 对草稿的评论。

E.F. 感谢高等研究院（Institute for Advanced Study）在工作期间多次访问时给予的款待。

% [原文疑误: 原书 p.13 此行将 Institute 误拼为 "Intitute"（"the Intitute for Advanced Study"）；译文按正确拼写译为「高等研究院」，不改变致谢内容。]
